\section{A substitution-stable single introduction}
\label{sec:diagonal-core}

The construction retains every numerical term and every formula of the
native scalar calculus.  It also retains the ordinary readings of its two
binders.  The change is confined to the introduction rule: a particular
\emph{syntactic diagonal shape} selects progress; every other matrix selects
cases.  The diagonal is chosen so that its zero instance is a theorem and
its recognition commutes with capture-free substitution.

\subsection{Syntax, readings, and all seven primitive rules}

Numerical terms and native formulas are
\[
 t ::= x\mid 0\mid S(t)\mid I(n,a;u.h),
 \qquad
 A ::= (s=t)\mid Q_y A\mid D_{x,y}A.
\]
In an iterator only $u$ is bound, and only in $h$; in $D$ the variables
$x,y$ are distinct.  We work up to renaming of bound variables, and all
substitutions avoid capture.  No restriction is imposed on public formulas,
intermediate formulas, iterator bodies, or witnesses.  The designated
contradiction remains $0=S0$.  The ordinary readings remain
\[
 Q_y B := B(0)\to\forall y\,B(y),\qquad
 D_{x,y}B := \exists x\bigl(B(x,0)\to\forall y\,B(x,y)\bigr).
\]

\begin{definition}[The diagonal predicate]
For a matrix $B(x,y)$, let $\operatorname{Diag}_y(B)$ mean that its
\emph{literal syntax} is
\begin{equation}
 B(x,y)=
 \bigl(I(y,y;u.h)=I(y,y;v.k)\bigr)
 \label{eq:diag-shape}
\end{equation}
for some bodies $h,k$.  The displayed counter and initializer on both sides
are each the variable bound as $y$ by the impending $D$.
The state binders $u,v$ may be renamed independently.  The bodies may use
$x,y$, their own state variable, and any fixed parameters.
This is syntactic recognition, without simplification or an equality test
inside the arithmetic theory.  Put
\[
 \eta_B(z)=
 \begin{cases}
 z,&\operatorname{Diag}_y(B),\\
 0,&\text{otherwise}.
 \end{cases}
\]
\end{definition}

The calculus $\mathcal G$ has the following complete primitive inventory.
Contexts are finite reusable lists of native formulas.
\begin{align*}
 \frac{A\in\Gamma}{\Gamma\vdash A}
   &\quad\text{Hyp},\\[3pt]
 \Gamma\vdash I(0,a;u.h)=a
   &\quad\text{I0},\\[3pt]
 \Gamma\vdash I(Sn,a;u.h)=h[I(n,a;u.h)/u]
   &\quad\text{IS},\\[3pt]
 \frac{\Gamma\vdash s=t\qquad
       \Gamma\vdash a(s)=b(s)}
      {\Gamma\vdash a(t)=b(t)}
   &\quad\text{Tr}_{=},\\[6pt]
 \frac{\Gamma\vdash Q_y B(y)\qquad\Gamma\vdash B(0)}
      {\Gamma\vdash B(t)}
   &\quad\text{QE},\\[6pt]
 \frac{\Gamma\vdash D_{x,y}B(x,y)\qquad
       \Gamma,Q_yB(b,y)\vdash A}
      {\Gamma\vdash A}
   &\quad\text{DE},\\[6pt]
 \frac{\Gamma,B(t,\eta_B(z))\vdash B(t,Sz)}
      {\Gamma\vdash D_{x,y}B(x,y)}
   &\quad\text{DI}_{\mathcal G}.
\end{align*}
In $\mathrm{Tr}_{=}$, $a(v),b(v)$ are numerical term templates; transport
inside a compound formula is not primitive.  In DE, $b$ is fresh for
$\Gamma,A$ and the fixed parameters of $B$.  In DI, $z$ is fresh for
$\Gamma,t$ and the fixed parameters of $B$.  The displayed local
assumptions are discharged.  There is no primitive reflexivity,
substitution, cut, ordinary connective, case, or additional induction
rule.

For comparison, $\mathcal N$ is the eight-rule native kernel of the
repository: its first six rules are the ones displayed above, and its two
introductions are
\begin{equation}
 \frac{\Gamma,B(t,0)\vdash B(t,Sz)}
      {\Gamma\vdash D_{x,y}B(x,y)}
 \quad\text{DI}_{\rm case},
 \qquad
 \frac{\Gamma,E(z)\vdash E(Sz)}
      {\Gamma\vdash Q_y E(y)}
 \quad\text{QI}_{\rm eq},
 \label{eq:diag-native-intros}
\end{equation}
where $E$ is a literal equation and the same freshness conventions apply.

\subsection{Why the diagonal survives substitution}

\begin{lemma}[Preservation and reflection of the guard]
\label{lem:diag-stable}
Let $\sigma$ be a capture-free substitution for the free parameters of
$B(x,y)$, with $x,y$ retained as the distinguished bound coordinates.
Then
\[
 \operatorname{Diag}_y(B\sigma)
 \quad\Longleftrightarrow\quad
 \operatorname{Diag}_y(B).
\]
The same assertion holds when a witness is substituted for $x$, provided
that the witness does not contain the bound coordinate $y$.
\end{lemma}
\begin{proof}
Rename binders away from every term in the substitution range.
A constructor-headed term remains headed by that constructor after
substitution.  Moreover, a term becomes the distinguished bound variable
$y$ if and only if it was already $y$: another retained variable stays
different from $y$, a substituted parameter has an image not containing
$y$, and a constructor-headed term cannot become a variable.
Consequently an iterator becomes headed by the exact pair of arguments
$(y,y)$ if and only if it had that pair before substitution.
The root of a formula is also preserved.  An entire side that is a
substituted parameter may become iterator-headed, but cannot acquire the
specified shape, because its image has no free occurrence of the protected
$y$.  An existing iterator matches the shape exactly when its two displayed
arguments match $y$, by the preceding argument.  Applying these facts to
both sides proves both directions of the equivalence.  The bodies are
unrestricted, so they impose no further condition.  Witness substitution
is the same argument with $x$ among the substituted parameters.
\end{proof}

The reflection direction is essential.  The superficially similar guard
\[
 I(y,0;u.h)=I(y,0;v.k)
\]
does not have it.  For a free parameter $p$, the matrix
\[
 I(y,p;u.0)=I(y,0;v.0)
\]
does not match the zero-initializer guard, whereas its substitution
$p:=0$ does.  Thus that test can change a case instruction into a progress
instruction.  Requiring the initializer itself to be the protected bound
variable $y$ prevents this event.

\begin{lemma}[Structural admissibility]
\label{lem:diag-structural}
Weakening, exchange, contraction, renaming, capture-free numerical
substitution, and ordinary-assumption cut are admissible in $\mathcal G$.
Numerical substitution can be performed by induction on a derivation
without changing any last-rule constructor or either mode of
$\mathrm{DI}_{\mathcal G}$.
\end{lemma}
\begin{proof}
Weakening, exchange and contraction change only hypothesis membership;
rename eigenparameters before inserting a context that could contain
them.  For numerical substitution, induct on the derivation.  The
computation and transport rules commute with capture-free substitution.
For QE and DE use the usual interchange of lifted substitution and
instantiation, after renaming their bound and eigenvariables.  In DI,
Lemma~\ref{lem:diag-stable} says that the same one of $0,z$ is retained
after substitution.  The substituted premise is therefore exactly the
premise of the same mode of DI.  This proves the induction step directly.

For cut, replace each use of a discharged ordinary hypothesis by its
supplied derivation.  Induct on the derivation in which that hypothesis
occurs, weakening the inserted proof in each local context.  Rename local
eigenparameters away from the inserted proof.  The guard is independent
of context and proof terms, so its mode is unchanged.  This proves both
single-assumption cut and simultaneous replacement of hypotheses by
derivations.  None of these arguments appeals to the equivalence with
$\mathcal N$ proved below.
\end{proof}

\subsection{Recovering every case introduction first}

\begin{lemma}[Elementary equality and the diagonal base]
\label{lem:diag-base}
Reflexivity, symmetry, transitivity and congruence for numerical terms are
derivable using only I0 and $\mathrm{Tr}_{=}$.  If
$\operatorname{Diag}_y(B)$, then $\Gamma\vdash B(t,0)$ for every
context $\Gamma$ and witness $t$.
\end{lemma}
\begin{proof}
Put $r=I(0,a;u.0)$.  The computation $r=a$, used as both transport
premises with equation template $v=a$, gives $a=a$.
From $s=t$ and $s=s$, the template $v=s$ gives $t=s$.
The template $s=v$ gives transitivity, and the template
$h(s)=h(v)$ gives congruence after reflexivity.
At $y=0$, both sides of~\eqref{eq:diag-shape} are of the form
$I(0,0;u.h')$, hence each equals $0$ by I0.
Symmetry and transitivity give their equality.  No induction, body
extensionality, or test for numerical equality is used.
\end{proof}

\begin{lemma}[Unrestricted case introduction]
\label{lem:diag-all-cases}
Every instance of $\mathrm{DI}_{\rm case}$ in
\eqref{eq:diag-native-intros} is admissible in $\mathcal G$, for every
native matrix $B$.
\end{lemma}
\begin{proof}
If $B$ does not match the diagonal, the desired rule is primitive.
Otherwise start with
$\Gamma,B(t,0)\vdash B(t,Sz)$.
By Lemma~\ref{lem:diag-base} the discharged base is already a theorem,
so cut it out.  This gives $\Gamma\vdash B(t,Sz)$.
Weaken by $B(t,z)$ and apply the progress mode of
$\mathrm{DI}_{\mathcal G}$.
This recovers the omitted case instance without any general equation
induction.
\end{proof}

\begin{lemma}[Conditional closure and ordinary numerical cases]
\label{lem:diag-case-interface}
For every native formula template $F(y)$, the rules
\[
 \frac{\Gamma,F(0)\vdash F(Sz)}{\Gamma\vdash Q_yF(y)},
 \qquad
 \frac{\Gamma\vdash F(0)\qquad\Gamma,F(0)\vdash F(Sz)}
      {\Gamma\vdash F(t)}
\]
are admissible in $\mathcal G$, with $z$ fresh for the retained
context and fixed parameters.  In particular, the second successor
premise may ignore its additional hypothesis $F(0)$.
\end{lemma}
\begin{proof}
Choose an unused witness coordinate $x$.  The recovered case rule at
witness $0$ introduces $D_{x,y}F(y)$ from the first displayed premise.
DE opens $Q_yF(y)$, which is independent of the fresh witness; return
that hypothesis as the conclusion of DE.  This proves the first rule.
QE with the base proof proves the second rule.
\end{proof}

\subsection{A small logical tool with no hidden induction}

Only implication between two equations is needed in the completeness
proof.  It can be constructed directly, so no general formula selector,
reflected equality, predecessor, or pairing theorem is needed at this
stage.  Define the numerical selector
\[
 C(w,a,b):=I(w,a;u.b),\qquad u\notin\operatorname{FV}(b).
\]
Computation gives $C(0,a,b)=a$ and $C(Sw,a,b)=b$.
For equations $U=(a=b)$ and $V=(c=d)$, with $w$ fresh, put
\begin{equation}
 U\Rightarrow_{=}V
   :=Q_w\bigl(C(w,a,c)=C(w,b,d)\bigr).
 \label{eq:diag-equational-implication}
\end{equation}
This is a native formula abbreviation, not an additional connective.

\begin{lemma}[Implication for equations]
\label{lem:diag-eq-implication}
For literal equations $U,V$, the rules
\[
 \frac{\Gamma,U\vdash V}{\Gamma\vdash U\Rightarrow_{=}V},
 \qquad
 \frac{\Gamma\vdash U\Rightarrow_{=}V\qquad\Gamma\vdash U}
      {\Gamma\vdash V}
\]
are derivable using computation, equational transport, and the
case-derived closure rule of Lemma~\ref{lem:diag-case-interface}.
\end{lemma}
\begin{proof}
Let $J(w)$ denote the equation inside the $Q$ in
\eqref{eq:diag-equational-implication}.
Its zero instance is interderivable with $U$, and each successor
instance is interderivable with $V$, by the two selector computations
and elementary equality reasoning.
For introduction, convert the retained $J(0)$ to $U$, insert the given
derivation of $V$, convert $V$ to $J(Sz)$, and apply case-derived
$Q$ introduction.  For elimination, convert the proof of $U$ to
$J(0)$, apply QE at $S0$, and convert $J(S0)$ to $V$.
All conversions are between equations.  This argument uses neither
compound transport nor general progress.
\end{proof}

\begin{lemma}[Cases in a context that mentions the inspected term]
\label{lem:diag-dependent-cases}
Let $E$ be an equation and let $r$ be any term.  The rule
\[
 \frac{\Gamma,r=0\vdash E\qquad\Gamma,r=Sv\vdash E}
      {\Gamma\vdash E}
\]
is admissible when $v$ is fresh for $\Gamma,r,E$.
The context $\Gamma$ may contain $r$ and its free variables.
\end{lemma}
\begin{proof}
Keep $\Gamma,r,E$ fixed and choose a new variable $y$.
Set $F(y):=((r=y)\Rightarrow_{=}E)$.
Discharge the two branch equations by
Lemma~\ref{lem:diag-eq-implication}; this gives proofs of $F(0)$ and
$F(Sv)$ in $\Gamma$.  Apply the ordinary case rule of
Lemma~\ref{lem:diag-case-interface} to $F(y)$ at target $r$.
The result is $(r=r)\Rightarrow_{=}E$.
Reflexivity and implication elimination give $E$.
The variable being generalized in the case inference was the fresh
$y$, not a variable occurring in $\Gamma$.  Thus the proof does not
use the invalid convention of changing a context silently in each branch.
\end{proof}

\subsection{Recovering unrestricted equation progress}

\begin{theorem}[The full native progress rule]
\label{thm:diag-progress}
For every literal equation $E(y)$, the rule
\[
 \frac{\Gamma,E(z)\vdash E(Sz)}{\Gamma\vdash Q_yE(y)}
\]
is admissible in $\mathcal G$, with the usual eigenparameter condition.
\end{theorem}
\begin{proof}
Write $E(y)=(s(y)=t(y))$ and let $\pi$ be the supplied step derivation.
Define the diagonal lift
\begin{equation}
 H(y):=
 \bigl(I(y,y;u.s(y))=I(y,y;u.t(y))\bigr),
 \qquad u\notin\operatorname{FV}(s,t).
 \label{eq:diag-lift}
\end{equation}
By computation and equality alone,
\begin{equation}
 \vdash H(0),\qquad
 H(Sv)\dashv\vdash E(Sv)
 \label{eq:diag-lift-endpoints}
\end{equation}
for every term $v$.  The second statement uses constant state updates;
it does not assert any identity between arbitrary iterator bodies.

Work temporarily in the retained context
$\Delta:=\Gamma,E(0)$, which contains no fresh induction variable.
For fresh $y$ set
\[
 \Phi(y):=H(y)\Rightarrow_{=}E(Sy).
\]
There are two case proofs in this fixed context:
\begin{itemize}
\item At $y=0$, substitute $0$ for the eigenparameter of $\pi$.
  The retained $E(0)$ gives $E(S0)$.
  The assumption $H(0)$ may be ignored, so equation-implication
  introduction gives $\Delta\vdash\Phi(0)$.
\item At $y=Sv$, assume $H(Sv)$.
  By~\eqref{eq:diag-lift-endpoints} it gives $E(Sv)$.
  Substitute $Sv$ for the eigenparameter of $\pi$ and obtain $E(SSv)$.
  Discharging $H(Sv)$ gives $\Delta\vdash\Phi(Sv)$.
\end{itemize}
These uses of substitution are justified directly by
Lemma~\ref{lem:diag-structural}; they do not invoke the result currently
being proved.
The ordinary case rule applies to the \emph{compound} native formula
$\Phi$ and yields $\Delta\vdash\Phi(z)$ for a fresh $z$.
Now assume $H(z)$, use equation-implication elimination, and convert
$E(Sz)$ back to $H(Sz)$.  We have proved
\[
 \Delta,H(z)\vdash H(Sz).
\]
The case inference was performed \emph{before} adding $H(z)$ to the
context, so its freshness condition is satisfied.  Equivalently, one
could use the explicitly justified contextual case rule of
Lemma~\ref{lem:diag-dependent-cases}.

Take an unused witness coordinate $x$ in the matrix $H(y)$.
It has the exact diagonal shape, so primitive progress DI at witness
$0$ gives $\Delta\vdash D_{x,y}H(y)$.
DE, with its fresh witness unused, gives $\Delta\vdash Q_yH(y)$.
Together with $H(0)$, QE at a fresh successor $Sz$ yields $H(Sz)$,
hence $E(Sz)$.  We have therefore constructed
\[
 \Gamma,E(0)\vdash E(Sz).
\]
Finally apply the already recovered case-derived $Q$ introduction to
obtain $\Gamma\vdash Q_yE(y)$.

The temporary retention and final discharge of $E(0)$ are deliberate.
No inference identifies conditional closure with unconditional progress.
In particular, the argument also works when the base equation is false.
\end{proof}

\subsection{Exact equivalence and the arithmetic consequence}

\begin{theorem}[Identity-language equivalence]
\label{thm:diag-equivalence}
For every finite context $\Gamma$ and every native formula $A$,
\[
 \boxed{\quad
 \Gamma\vdash_{\mathcal G}A
 \quad\Longleftrightarrow\quad
 \Gamma\vdash_{\mathcal N}A.
 \quad}
\]
Both translations act effectively on finite derivations, and neither
translation changes a term, formula, or retained assumption.
\end{theorem}
\begin{proof}
Translate a $\mathcal G$ derivation to $\mathcal N$ by induction.
The first six rules are shared.  An unguarded DI instance is native
case introduction.  In a guarded DI instance the matrix is an equation.
The native progress rule applied to $E(y):=B(t,y)$ gives $Q_yB(t,y)$.
Under the temporary base $B(t,0)$, QE supplies $B(t,Sz)$; native case
DI then introduces the required $D_{x,y}B(x,y)$.

Conversely, translate a native derivation recursively.
The shared rules remain unchanged, every native case introduction is
replaced by Lemma~\ref{lem:diag-all-cases}, and every native equation
progress introduction is replaced by
Theorem~\ref{thm:diag-progress}.  All constructions use the original
retained context and satisfy the freshness conditions.
\end{proof}

The proof has a strict dependency order:
guard stability; structural admissibility; elementary equality;
all case introductions; case-derived closure and numerical cases;
implication between equations; general equation progress.
It does not use equality reflection, a scalar predecessor, pairing,
primitive-recursion compilation, or the eventual bi-interpretation in
order to establish equivalence.

\begin{proposition}[Transfer of the established bi-interpretation]
\label{prop:diag-ipra}
Compose Theorem~\ref{thm:diag-equivalence} with the native scalar
calculus's established bi-interpretation with iPRA.  This gives a
bi-interpretation of $\mathcal G$ with iPRA, and equiconsistency.
The numerical domains and interpreted equality are unchanged, and both
numerical composites have the identity graph $x=y$.
\end{proposition}
\begin{proof}
The first component of each composite, between $\mathcal G$ and
$\mathcal N$, is literally the identity on the full syntax and on all
sequents.  It therefore preserves the established whole-domain
interpretations, their internal formula round trips, and their numerical
identity graphs.  The equivalence also preserves the designated native
contradiction $0=S0$ in the empty context.  Composing the finite proof
translations in each direction preserves, and reflects, derivability of
contradiction, which proves equiconsistency.
\end{proof}

\subsection{A sufficient criterion for other guards}

\begin{proposition}[The guard criterion]
\label{prop:diag-guard-criterion}
Let $\mathsf P_y(B)$ be a decidable predicate on native matrices that
respects capture-free renaming.  Suppose that:
\begin{enumerate}
\item Every matrix satisfying $\mathsf P$ is a literal equation, and
  its instance $B(t,0)$ is derivable for every witness $t$ using only
  computation and equational transport.
\item Every equation $E(y)$ has an equation lift $H(y)$, independent
  of an added unused witness coordinate $x$, such that
  $\mathsf P_y(H)$ holds and
  $H(Sv)\dashv\vdash E(Sv)$ is derivable using only computation
  and equational transport, for arbitrary $v$.
\end{enumerate}
Replace the diagonal test in the seven-rule presentation by $\mathsf P$.
The resulting calculus $\mathcal G_{\mathsf P}$ derives exactly the
native sequents.  If, in addition, $\mathsf P$ is preserved and reflected
by capture-free numerical substitution, its structural substitution
theorem follows by direct rule-preserving induction.
\end{proposition}
\begin{proof}
The first condition recovers every missing case introduction exactly as
in Lemma~\ref{lem:diag-all-cases}.  Hence the case-derived closure,
numerical cases and equation-implication rules are available.
The second condition supplies the only properties of the lift used in
Theorem~\ref{thm:diag-progress}.  Since the base and endpoint conversions
belong to the fragment with computation and equational transport, neither
condition presupposes the general progress rule being recovered.

For clarity, equivalence does not silently assume numerical substitution
in $\mathcal G_{\mathsf P}$ when the optional stability condition is
absent.  Translate native derivations by strong induction on their height.
In the native progress case, substitute $0$ and $Sv$ into its native
premise derivation first.  Native substitution preserves derivation
height, so both resulting proofs still have smaller height than the
progress conclusion and can be translated by the induction hypothesis.
They supply the two $\Phi$ branches in the displayed proof above.
Renaming, weakening and cut require only the stated renaming invariance.
The other native rules translate as before, and the reverse translation
uses only the fact that a progress-selected matrix is an equation.
This proves the sequent equivalence.  Under the additional preservation
and reflection condition, numerical substitution instead follows directly
from the proof of Lemma~\ref{lem:diag-structural}, with no change of mode.
\end{proof}

\subsection{What is simpler, and what is not being claimed}

Relative to the eight-rule native kernel, this construction removes one
primitive introduction while retaining its full language, all witnesses,
and both original binder readings.  Relative to the earlier seven-rule
witness-selected presentation, its improvement is structural: the mode
depends only on protected bound-variable positions in the matrix, so
ordinary substitution is proved by a direct rule-preserving induction.
There is no exceptional substitution-repair proof.

The earlier guarded-language proposal restricted the public atomic
matrices; here every original formula remains available and the comparison
theorem is the identity on \emph{all} sequents.  The earlier
uniform-progress proposal changes the meaning of $Q$ at equations and
modifies its elimination; here neither change is made.  The diagonal
initializer supplies the missing combination of a provable zero base and
a substitution-invariant test.

This is not a proof of global minimality.  A seven-rule count is not a
new record by itself, and the syntactic guard has an implementation cost.
Any claim about a smaller complete checker must be supported by the
separate matched source measurements, including syntax validation,
substitution, freshness, and the public certificate interface.
The new equivalence above is a written mathematical proof; it is not, by
itself, a claim that a Lean or Rocq verification of this new kernel has
been compiled.
